% TEMPLATE-MILC-v3.tex - plantilla maestra UNICA de las guias MILC v3.
%
% La usan las cuatro familias de guias, cada una con su driver:
%   · Grados 6-11          -> scripts/build-guias-g11.py
%   · Bebras               -> scripts/build-guias-bebras.py
%   · Semillero            -> scripts/build-guias-semillero.py
%   · Territorio interior  -> scripts/build-guias-territorio-interior.py
%
% Antes habia tres copias casi identicas (~400 lineas iguales, 6 distintas):
% cada mejora habia que aplicarla tres veces, y en la practica no se hacia
% (el motor de recursos vivia solo en dos de ellas). Lo que varia entre
% familias esta parametrizado en 6 placeholders que cada driver llena
% (se nombran aqui SIN su sintaxis real: el builder reemplaza por texto plano,
% tambien dentro de los comentarios, y un valor multilinea rompe el preambulo):
%
%   HEADER_IZQ        encabezado izquierdo de pagina
%   HEADER_DER        encabezado derecho de pagina
%   PORTADA_ETIQUETA  rotulo sobre el numero grande (GRADO/MOMENTO/MODULO)
%   PORTADA_SERIE     linea de serie de la portada
%   PORTADA_ID        identificador de la guia en la portada
%   PORTADA_CIRCULO   el circulito de grado (°), o vacio si no aplica
%   PDF_TITULO        titulo en texto plano (sin \textbf ni comillas TeX) para
%                     los metadatos del PDF (/Title); lo produce lib_xelatex.texto_pdf
%
% Composicion (auditoria 2026-09-03): las bandas de estacion piden espacio
% (\Needspace*) y prohiben el salto justo despues (\nopagebreak), asi no quedan
% huerfanas al pie; las cajas largas (softbox, drawbox, infoband, puentebox) son
% `breakable`, asi una caja de media pagina ya no arrastra todo a la siguiente
% dejando la anterior al 40 %. Todo texto sobre mostaza o verde va en milcNegro
% (blanco sobre #E5B400 da 1,93:1 y sobre #58A923 2,95:1; ninguno pasa WCAG AA).
% El resto de claves las documenta content/guias/_SCHEMA.md. El script valida
% que no queden placeholders sin reemplazar antes de escribir el archivo final.

% Etiquetado PDF (latex-lab): /Lang, estructura y texto alternativo de las
% figuras, para lectores de pantalla. Debe ir ANTES de \documentclass.
\DocumentMetadata{lang=es-CO,tagging=on}
\documentclass[11pt,letterpaper]{article}
% headheight=24pt: la cabecera derecha lleva dos lineas (identidad de la guia
% y estacion actual). headsep se acorta en la misma medida para que el bloque
% de texto quede exactamente donde estaba con headheight=12pt.
\usepackage[margin=2.0cm,headheight=24pt,headsep=13pt]{geometry}
\usepackage[table]{xcolor}
\usepackage{fontspec}
\usepackage[spanish]{babel}
\usepackage{tikz}
% Librerías TikZ para los diagramas de las guías (bloque `recursos:`):
%  · babel  → babel en español vuelve activos caracteres como «>», y TikZ los usa
%             en sus opciones (>=stealth, ->). Sin esta librería, cualquier
%             diagrama con flechas revienta con «\language@active@arg> has an
%             extra }». Debe ir ANTES de que se dibuje nada.
%  · positioning        → sintaxis «below left=2pt and 3pt».
%  · decorations.pathreplacing → llaves ({brace}) para señalar rangos.
\usetikzlibrary{babel,positioning,decorations.pathreplacing}
\usepackage{graphicx}
% Circuitos: el trabajo de electrónica se hace hoy con micro:bit y sus sensores
% integrados (MakeCode, sin cableado), así que no se necesitan esquemáticos.
% Si algún día entran montajes con protoboard, basta descomentar esta línea:
% los diagramas .tex/.tikz declarados en `recursos:` ya se incrustan solos.
% \usepackage{circuitikz}
\usepackage[most]{tcolorbox}
\usepackage{tabularx}
\usepackage{array}
\usepackage{fancyhdr}
\usepackage{titlesec}
\usepackage[document]{ragged2e}  % bandera a la derecha en todo el cuerpo
\usepackage{enumitem}
\usepackage{microtype}
\usepackage{etoolbox}
\usepackage{needspace}
% hyperref al final del preambulo: metadatos del PDF (titulo, autor, idioma,
% familia), marcadores por estacion (\pdfbookmark) y enlaces sin recuadro.
\usepackage[hidelinks,
  pdftitle={Identidad prestada: lo que se compra, lo que se hace y lo que se paga},
  pdfauthor={PhD. Álvaro Cárdenas Orozco},
  pdfsubject={Territorio interior · Educación socioemocional · Grado 8° · Momento 6 · MILC v3},
  pdflang={es-CO},
  bookmarksopen=true,bookmarksnumbered=false]{hyperref}

% Helvetica Neue no trae la flecha U+2192, asi que cada «→» escrito literalmente
% en el YAML se compilaba a NADA, en silencio (462 flechas en 61 guias: rutas del
% tipo «paleta → tipografia → lockup» salian rotas). Mapeamos el caracter a su
% version matematica, que si tiene glifo. Asi el autor puede seguir escribiendo
% «→» comodamente en el YAML.
\usepackage{newunicodechar}
\newunicodechar{→}{\ensuremath{\rightarrow}}
% Mismo problema con los simbolos de marca y vinetas: el «✓» de «una palomita ✓»
% salia como cuadrito vacio en el PDF de todas las guias que lo usaban.
\usepackage{pifont}
\newunicodechar{✓}{\ding{51}}
\newunicodechar{✗}{\ding{55}}
\newunicodechar{✔}{\ding{52}}
\newunicodechar{•}{\textbullet}
\newunicodechar{≥}{\ensuremath{\geq}}
\newunicodechar{≤}{\ensuremath{\leq}}

\setmainfont{Helvetica Neue}
\setsansfont{Helvetica Neue}
\setlength{\parindent}{0pt}
\setlength{\parskip}{6pt}
% Sin particion de palabras: en 6.º-7.º una palabra cortada por guion al final
% de linea («Comuni-cacion») frena la lectura mucho mas que una linea corta.
\hyphenpenalty=10000
\exhyphenpenalty=10000
\pretolerance=2000
\tolerance=3000
\setlength{\emergencystretch}{3em}
\linespread{1.10}
\renewcommand{\arraystretch}{1.25}
\setlist{leftmargin=5mm,itemsep=1mm,topsep=1mm}
\newcolumntype{Y}{>{\RaggedRight\arraybackslash}X}

\definecolor{milcMagenta}{HTML}{E6007E}
\definecolor{milcVino}{HTML}{5A0038}
\definecolor{milcTurquesa}{HTML}{0093A5}
\definecolor{milcVerde}{HTML}{58A923}
\definecolor{milcMostaza}{HTML}{E5B400}
\definecolor{milcOcre}{HTML}{B86600}
\definecolor{milcNegro}{HTML}{241816}
\definecolor{milcGris}{HTML}{F7F7F4}
\definecolor{milcLinea}{HTML}{E8E2DD}
\definecolor{milcGradePrimary}{HTML}{0D9488}
\definecolor{milcGradeDark}{HTML}{115E59}
\definecolor{milcGradeSoft}{HTML}{CCFBF1}
\definecolor{milcGradeAccent}{HTML}{CFE7FF}
\definecolor{milcGradeText}{HTML}{FFFFFF}
\color{milcNegro}

\pagestyle{fancy}
\fancyhf{}
% Cabecera de dos lineas: identidad de la guia arriba y, a la derecha y en
% gris, la estacion en curso (\markright desde \milcestacion). Antes de la
% primera estacion la segunda linea va vacia; el \strut mantiene la altura
% para que la linea de arriba no baile entre paginas.
\lhead{\small Territorio interior · Educación socioemocional\\[-1pt]{\footnotesize\strut}}
\rhead{\small Grado 8° · Momento 6 · MILC v3\\[-1pt]{\footnotesize\strut\color{milcNegro!65}\rightmark}}
\cfoot{\small\thepage}
\renewcommand{\headrulewidth}{0pt}

% Sobre mostaza (#E5B400) y verde (#58A923) el texto va en milcNegro; sobre el
% resto de la paleta, en blanco. Se usa dentro de fonttitle/fontupper, donde un
% \color posterior manda sobre coltitle.
\newcommand{\milctintasobre}[1]{%
  \ifboolexpr{test{\ifstrequal{#1}{milcMostaza}} or test{\ifstrequal{#1}{milcVerde}}}%
    {\color{milcNegro}}{\color{white}}}

\newtcolorbox{titlebox}[1]{
  enhanced,colback=#1,colframe=#1,arc=7mm,boxrule=0pt,
  left=7mm,right=7mm,top=3mm,bottom=3mm,
  before skip=8pt,after skip=8pt,
  fontupper=\sffamily\bfseries\Large\color{white}
}

\newtcolorbox{softbox}[3]{
  enhanced,breakable,pad at break=3mm,
  colback=#2,colframe=#1,borderline west={4pt}{0pt}{#1},
  arc=2mm,boxrule=.4pt,left=4mm,right=4mm,top=3mm,bottom=3mm,
  before skip=6pt,after skip=8pt,title={#3},coltitle=white,
  colbacktitle=#1,fonttitle=\sffamily\bfseries\milctintasobre{#1},fontupper=\RaggedRight,
  attach boxed title to top left={xshift=3mm,yshift=-2mm},
  boxed title style={arc=2mm, boxrule=0pt}
}

\newtcolorbox{drawbox}[2]{
  enhanced,breakable,pad at break=4mm,
  colback=white,colframe=#1,arc=4mm,boxrule=1pt,
  left=5mm,right=5mm,top=4mm,bottom=4mm,before skip=8pt,after skip=8pt,
  title={#2},coltitle=white,colbacktitle=#1,fonttitle=\sffamily\bfseries\milctintasobre{#1},
  fontupper=\RaggedRight,
  attach boxed title to top left={xshift=4mm,yshift=-2mm},
  boxed title style={arc=3mm, boxrule=0pt}
}

\newtcolorbox{infoband}[1]{
  enhanced,breakable,pad at break=2mm,colback=#1!8,colframe=#1!50,arc=2mm,boxrule=.5pt,
  left=4mm,right=4mm,top=2mm,bottom=2mm,before skip=4pt,after skip=4pt,
  fontupper=\small\RaggedRight
}

% Puente narrativo entre secciones (contrato editorial Conéctate).
% Da continuidad: "Acabas de X. Ahora vas a Y porque Z."
% Visualmente: banda izquierda en color del grado, texto en itálica pequeña.
\newtcolorbox{puentebox}{
  enhanced,breakable,pad at break=2mm,colback=milcGradeSoft,colframe=milcGradeDark,
  borderline west={3pt}{0pt}{milcGradeDark},
  arc=0mm,boxrule=0pt,left=5mm,right=4mm,top=2.5mm,bottom=2.5mm,
  before skip=4pt,after skip=8pt,
  fontupper=\itshape\small\color{milcNegro}\RaggedRight
}
% La flecha va en modo matematico a proposito: Helvetica Neue NO tiene el glifo
% U+2192, asi que \textrightarrow se compilaba a NADA («Missing character: There
% is no → in font Helvetica Neue») y el puente salia sin flecha. En modo
% matematico la flecha viene de la fuente matematica y si se dibuja.
\newcommand{\puente}[1]{%
  \begin{puentebox}%
    \textcolor{milcGradeDark}{$\rightarrow$}\hspace{2mm}#1%
  \end{puentebox}%
}

\newcommand{\linead}[1]{\noindent\color{milcLinea}\rule{#1}{0.5pt}\color{milcNegro}}
\newcommand{\respuesta}[1]{\par\vspace{2mm}\foreach \n in {1,...,#1}{\noindent\color{milcLinea}\rule{\linewidth}{0.5pt}\par\vspace{5mm}}\color{milcNegro}}
\newcommand{\checkbox}{\textcolor{milcGradeDark}{\fbox{\phantom{x}}}\hspace{5pt}}

% ─── Listas aireadas para las guías socioemocionales ────────────────────
% Componer las verificaciones como «\checkbox texto\\» deja el texto que salta
% de línea debajo del cuadrito y apelmaza el bloque. Estos entornos alinean la
% continuación y airean los ítems. Los usa Territorio Interior; cualquier
% familia puede hacerlo.
\newlist{checklista}{itemize}{1}
\setlist[checklista]{
  label=\textcolor{milcGradeDark}{\fbox{\phantom{x}}},
  leftmargin=9mm, labelsep=3.2mm, itemsep=2.4mm,
  topsep=2.5mm, parsep=0pt, partopsep=0pt, before=\RaggedRight
}
\newlist{pasoslista}{enumerate}{1}
\setlist[pasoslista]{
  label=\textcolor{milcGradeDark}{\sffamily\bfseries\arabic*.},
  leftmargin=8mm, labelsep=3mm, itemsep=2.8mm,
  topsep=1mm, parsep=0pt, partopsep=0pt, before=\RaggedRight
}

\newcommand{\phasecard}[4]{
\begin{tcolorbox}[enhanced,colback=#3,colframe=#2,arc=3mm,boxrule=.5pt,height=3.05cm,valign=center,halign=center,left=1.5mm,right=1.5mm,top=2mm,bottom=2mm]
{\sffamily\bfseries\fontsize{22}{24}\selectfont\textcolor{#2}{#1}}\par
{\sffamily\bfseries\small #4}
\end{tcolorbox}
}

% ── Piezas del rediseno 2026 (legibilidad 6.º-11.º) ───────────────────
% Banda de estacion: reemplaza el titlebox macizo. Titulo a la izquierda,
% dato practico (tiempo/modalidad) a la derecha, en una sola linea.
% Nunca huerfana: pide 9 lineas libres (o salta de pagina rellenando con
% \vfill) y prohibe el salto justo despues de la banda. Nueve y no seis:
% la caja partible que sigue necesita ~80pt (titulo adosado, relleno y dos
% lineas) para arrancar en la misma pagina; con menos, tcolorbox la manda
% entera a la siguiente y la banda se queda sola. El penalty 10000 va
% en `after`, ANTES del espacio posterior: un `after skip` (aunque sea 0pt)
% es un \vskip tras la caja, y ese glue es un punto de corte legal que un
% \nopagebreak escrito despues de \end{tcolorbox} ya no cierra. Cada banda
% deja marcador PDF y actualiza la cabecera derecha.
\newcounter{milcestacion}
\newcommand{\milcestacion}[2]{%
\Needspace*{9\baselineskip}%
\stepcounter{milcestacion}%
\pdfbookmark[1]{#2}{milcestacion\themilcestacion}%
\markright{#2}%
\begin{tcolorbox}[enhanced,colback=#1,colframe=#1,arc=2mm,boxrule=0pt,
  left=5mm,right=5mm,top=2.6mm,bottom=2.6mm,before skip=11pt,
  after={\par\nopagebreak\vskip7pt}]
{\sffamily\bfseries\fontsize{15}{18}\selectfont\milctintasobre{#1}#2}
\end{tcolorbox}}

% Tarjeta de paso numerada: sustituye las tablas «Pilar / Aplicacion», que
% eran formato de consulta docente, no de estudiante. El numero va en un
% circulo sobre el borde para que la secuencia se lea de un vistazo.
\newcommand{\milcpaso}[3]{%
\Needspace{4\baselineskip}%
\begin{tcolorbox}[enhanced,colback=white,colframe=milcLinea,arc=1.5mm,boxrule=.9pt,
  left=14mm,right=4mm,top=3mm,bottom=3mm,before skip=4pt,after skip=4pt,
  fontupper=\RaggedRight,
  overlay={\node[anchor=north west,circle,fill=#2,text=white,inner sep=0pt,
    minimum size=7.4mm,font=\sffamily\bfseries\small]
    at ([xshift=3.6mm,yshift=-3.2mm]frame.north west) {#1};}]
#3
\end{tcolorbox}}

% Casilla marcable de tamano util con lapiz (la anterior era \fbox{\phantom{x}},
% de alto variable segun el texto que la seguia).
\newcommand{\milcmarca}{\framebox[3.8mm]{\rule{0pt}{3.4mm}}}

% Fila marcable de lista suelta (checks de la estacion 1).
\newcommand{\milccheckrow}[1]{%
\par\vspace{1.8mm}\noindent
\makebox[6.4mm][l]{\raisebox{-.55mm}{\textcolor{milcNegro}{\milcmarca}}}%
\parbox[t]{\dimexpr\linewidth-6.4mm\relax}{\RaggedRight #1}\par}

% Celda marcable dentro de la rubrica (tabla de autoevaluacion). Misma
% sangria francesa, pero pensada para vivir dentro de una columna X.
\newcommand{\milcopcion}[1]{%
\makebox[5.2mm][l]{\raisebox{-.55mm}{\textcolor{milcNegro}{\milcmarca}}}%
\parbox[t]{\dimexpr\linewidth-5.2mm\relax}{\RaggedRight #1}}

% ── Recursos visuales de la guía ──────────────────────────────────────
% Declarados en el bloque `recursos:` del YAML y emitidos por el builder.
% \guiaFigura   → imagen rasterizada o PDF (foto, carta celeste, montaje).
% \guiaDiagrama → fuente TikZ/circuitikz (.tex) incrustada como vector:
%                 es la vía para circuitos y esquemáticos editables.
% [#1] = texto alternativo (alt) · #2 = ruta relativa al .tex · #3 = pie
% El alt llega al PDF etiquetado (/Alt) y lo lee el lector de pantalla.
\newcommand{\guiaFigura}[3][]{%
\begin{center}
\includegraphics[width=0.88\linewidth,height=0.38\textheight,keepaspectratio,alt={#1}]{#2}\\[1.5mm]
{\footnotesize\itshape\textcolor{milcNegro!75}{#3}}
\end{center}
}

\newcommand{\guiaDiagrama}[2]{%
\begin{center}
\input{#1}\\[1.5mm]
{\footnotesize\itshape\textcolor{milcNegro!75}{#2}}
\end{center}
}

\begin{document}
\thispagestyle{empty}

% ── Portada ───────────────────────────────────────────────────────────
% Antes: un rectangulo de color a pagina completa. Se veia bien en pantalla,
% pero en la fotocopiadora del colegio se comia el toner y salia gris sucio.
% Ahora la banda ocupa el tercio superior y el resto va en blanco.
\begin{tikzpicture}[remember picture,overlay]
  \path[fill=milcGradePrimary]
    (current page.north west) rectangle ([yshift=-10.6cm]current page.north east);

  \node[anchor=north west,text=milcGradeText] at ([xshift=2.0cm,yshift=-1.55cm]current page.north west)
    {\sffamily\bfseries\fontsize{12}{14}\selectfont MOMENTO};
  \node[anchor=north west,text=milcGradeText,opacity=.85] at ([xshift=2.0cm,yshift=-2.35cm]current page.north west)
    {\sffamily\bfseries\fontsize{8.5}{10}\selectfont TERRITORIO INTERIOR · SOCIOEMOCIONAL};
  \node[anchor=north east,text=milcGradeText,opacity=.85,align=right] at ([xshift=-2.0cm,yshift=-1.55cm]current page.north east)
    {\sffamily\bfseries\fontsize{9}{12}\selectfont I.E. Sor María Juliana\\Cartago, Valle del Cauca};

  \node[anchor=west,text=milcGradeText] (gradonum) at ([xshift=1.85cm,yshift=-7.1cm]current page.north west)
    {\sffamily\bfseries\fontsize{112}{112}\selectfont 6};
  % El circulito de grado (°) solo aplica a las guias de grado («11°»); en
  % Bebras y Semillero el numero grande es un momento/modulo, no un grado.
  % Se posiciona relativo al nodo `gradonum`, no en coordenadas absolutas.
  

  \node[anchor=west,text=milcGradeText,text width=11.4cm,align=left] at ([xshift=1.5cm,yshift=0.5cm]gradonum.east)
    {
     {\sffamily\bfseries\fontsize{9}{11}\selectfont\MakeUppercase{Territorio interior · Socioemocional}}\\[3mm]
     {\sffamily\bfseries\fontsize{21}{25}\selectfont\setlength{\baselineskip}{27pt}Identidad\\[4pt]prestada\\[4pt]comprar no es lo mismo que hacer}\\[3.5mm]
     {\sffamily\bfseries\fontsize{15}{18}\selectfont Grado 8° · Momento 6}
    };
\end{tikzpicture}

\vspace*{9.4cm}
\vfill

{\sffamily\bfseries\fontsize{8}{10}\selectfont\textcolor{milcGradeDark}{LO QUE TE LLEVAS HOY}}

\begin{tcolorbox}[enhanced,colback=white,colframe=milcNegro,arc=2mm,boxrule=1.6pt,
  left=5mm,right=5mm,top=4mm,bottom=4mm,before skip=3pt,after skip=12pt,
  fontupper=\RaggedRight\fontsize{11}{15.5}\selectfont]
Tu inventario de señales: lo que en tu curso «hay que tener» con su precio real y quién queda por fuera por no tenerlo, más una cosa que sepas hacer con las manos o con la cabeza y un plan de cuatro semanas para volverla visible.
\end{tcolorbox}

\begin{tcolorbox}[enhanced,colback=milcGris,colframe=milcGris,arc=2mm,boxrule=0pt,
  left=5mm,right=5mm,top=3.5mm,bottom=3.5mm,after skip=12pt]
{\sffamily\bfseries\fontsize{8}{10}\selectfont\textcolor{milcNegro!55}{ESTA GUÍA ES DE}}\\[3mm]
\begin{tabularx}{\linewidth}{X p{3.2cm} p{3.2cm}}
\linead{\linewidth} & \linead{3.0cm} & \linead{3.0cm}\\[-1mm]
{\fontsize{8.5}{10}\selectfont\textcolor{milcNegro!55}{Nombre completo}} &
{\fontsize{8.5}{10}\selectfont\textcolor{milcNegro!55}{Grupo}} &
{\fontsize{8.5}{10}\selectfont\textcolor{milcNegro!55}{Fecha}}\\
\end{tabularx}
\end{tcolorbox}

\vfill

\noindent\textcolor{milcLinea}{\rule{\linewidth}{1pt}}\\[2mm]
{\sffamily\bfseries\fontsize{9.5}{12}\selectfont Autor: Dr. Álvaro Cárdenas Orozco}\\[1mm]
{\sffamily\fontsize{8.5}{11}\selectfont\textcolor{milcNegro!55}{Metodología MILC · Apertura ancestral · Marca, moda y pertenencia · Triángulo Dussel-Estoicismo-Floridi}}

\newpage

\section*{Apertura: Identidad prestada: lo que se compra, lo que se hace y lo que se paga}

\begin{softbox}{milcGradeDark}{milcGradeSoft}{Saber ancestral}
En el Norte del Valle hay una identidad que \textbf{no se compró: se hizo}. \textbf{Cartago} es conocida en el país por el \textbf{calado} ---esa técnica de sacar hilos de la tela y bordar sobre la rejilla que queda---, y en \textbf{Ansermanuevo}, a veinte minutos, el calado y el bordado también se celebran como \textbf{fiesta propia del pueblo}. \emph{Piensa en lo raro y lo enorme que es eso:} \textbf{una región entera decidió que su signo de identidad fuera un trabajo hecho a mano, lento, que no se puede fabricar rápido y que lleva pegado el nombre de quien lo hizo}. \textbf{Y hay un detalle que las caladoras conocen bien:} una pieza de calado \textbf{se puede reconocer}. Quien sabe mirar distingue la mano, la puntada, el taller de donde salió. \emph{Es exactamente lo contrario de una marca:} \textbf{una marca la tienen millones de personas idénticas y no dice nada de ninguna}. \emph{Un calado lo tiene una sola, y dice quién lo hizo y cuánto tardó.} \textbf{Guarda esa diferencia, porque es de lo que trata toda esta guía:} \emph{hay señales que se compran y señales que se hacen, y no cuestan ni duran lo mismo.}
\end{softbox}

\begin{softbox}{milcTurquesa}{milcGris}{Saber contemporáneo}
¿Y qué pasa cuando alguien busca en las cosas lo que le falta en otro lado? \textbf{Se midió, y en grande.} \textbf{Helga Dittmar} y su equipo reunieron \textbf{753 resultados de 259 estudios} para responder si el \textbf{materialismo} ---darle mucha importancia a tener cosas y al dinero--- se relaciona con el bienestar. \textbf{La respuesta fue consistente: a más materialismo, \emph{menos} bienestar} (Dittmar, Bond, Hurst y Kasser, 2014). \textbf{Y lo más interesante no es el resultado sino la explicación.} Los análisis sugieren que el vínculo negativo se explica porque \textbf{las necesidades psicológicas quedan mal satisfechas}. \emph{Dicho simple:} uno compra buscando algo ---sentirse aceptado, sentirse capaz, sentirse alguien--- \textbf{y la compra no entrega eso}, porque no es algo que se venda. \textbf{Entonces uno vuelve a comprar.} \emph{Y encontraron algo más, que explica por qué esto pesa tanto en un salón:} el efecto es más fuerte \textbf{en contextos que apoyan los valores materialistas}. \textbf{Traducido:} \emph{no es solo cuánto te importe a ti la marca; es cuánto le importa al lugar donde estás metido ocho horas al día.}
\end{softbox}

\begin{softbox}{milcMostaza}{milcGris}{Pregunta-puente}
\textit{Un calado \textbf{se reconoce}: dice quién lo hizo. \textbf{¿Qué cosa de las que tienes te representa realmente a ti} \ldots \textbf{y cuál compraste para que no te dejaran por fuera}?}
\end{softbox}

\begin{softbox}{milcVerde}{milcGris}{Saber y saber hacer}
Al terminar podrás: (1) \textbf{distinguir} lo que compras por gusto de lo que compras como \textbf{señal} de pertenencia; (2) \textbf{explicar} por qué una señal comprada se devalúa y por qué comprar no satisface la necesidad que empujó la compra; (3) \textbf{identificar} a quién deja por fuera, en tu curso, la lista de lo que «hay que tener»; (4) \textbf{planear} cuatro semanas para volver visible algo que sabes hacer ---una señal que no se compra y no se devalúa---.
\end{softbox}



\small
\begin{tabularx}{\textwidth}{>{\bfseries}p{3.0cm}Y}
\rowcolor{milcGradePrimary!12}Autor & Dr. Álvaro Cárdenas Orozco\\
DBA & Comprende los fenómenos que afectan a su territorio, regula sus emociones ante ellos y acompaña a otros con empatía (transversal 6°--11°, articulado con la política «Escuela, territorio de vida» del MEN, Resolución 006519 de 2025, y la Ley 1620 de 2013)\\
Referentes & Primera ayuda psicológica (OMS, 2011/2012) · Directrices IASC (2007) · USGS y Servicio Geológico Colombiano · Pueblo nasa y la minga (CRIC) · Dussel · Estoicismo · Floridi\\
Producto final & Tu inventario de señales: lo que en tu curso «hay que tener» con su precio real y quién queda por fuera por no tenerlo, más una cosa que sepas hacer con las manos o con la cabeza y un plan de cuatro semanas para volverla visible.\\
\end{tabularx}
\normalsize

\puente{En el momento 1 hablamos del \textbf{peaje}. Hoy miramos el peaje que se paga \textbf{en plata} ---y el que lo paga la familia sin que nadie lo diga en voz alta---.}

\milcestacion{milcGradeDark}{Ruta MILC: así vamos a aprender}
\begin{tabularx}{\textwidth}{>{\centering\arraybackslash}X>{\centering\arraybackslash}X>{\centering\arraybackslash}X>{\centering\arraybackslash}X}
\phasecard{1}{milcVerde}{milcGris}{Escucha\\[-1mm]\small Observo, escucho y pregunto.} &
\phasecard{2}{milcTurquesa}{milcGris}{Sistematización\\[-1mm]\small Distingo y hago.} &
\phasecard{3}{milcMagenta}{milcGris}{Praxis\\[-1mm]\small Construyo y aplico.} &
\phasecard{4}{milcVino}{milcGris}{Evaluación\\[-1mm]\small Reflexión liberadora.}
\end{tabularx}


\puente{Primero, la lista honesta: qué «hay que tener» en tu curso. \emph{Todos la conocen y casi nunca se escribe.}}

\milcestacion{milcVerde}{Estación 1 de 3 · Escucha}

\begin{softbox}{milcVerde}{milcGris}{Escena para escuchar}
\textbf{Actividad 1 · IDENTIFICA --- Lo que «hay que tener».} \textbf{Nadie va a leer esta página.} \textbf{Primera parte.} Escribe la lista real de \textbf{lo que hay que tener en tu curso para que no te digan nada}: la marca de los tenis, el celular, el corte de pelo, la maleta, los audífonos, la ropa del día de la civil. \emph{No la lista oficial: la de verdad.} \textbf{Segunda parte.} Al frente de cada cosa, escribe \textbf{cuánto cuesta} ---búscalo si no lo sabes--- y \textbf{suma}. \emph{El número final es el punto de la actividad.} \textbf{Tercera parte.} Marca cuáles tienes tú y \textbf{por qué las compraste}: \emph{porque te gustaba} o \emph{porque no querías quedar por fuera}. \textbf{Sé honesto: casi siempre hay de las dos.} \textbf{Y la incómoda:} ¿alguna vez has pedido algo en tu casa \textbf{sabiendo que era difícil}? ¿Y sabes qué se dejó de hacer para comprarlo?
\end{softbox}

{\sffamily\bfseries\fontsize{8}{10}\selectfont\textcolor{milcNegro!55}{MARCA CUANDO LO HAYAS HECHO}}
\milccheckrow{Escribiste la lista real de lo que «hay que tener», no la oficial.}
\milccheckrow{Averiguaste los precios y \textbf{sumaste}.}
\milccheckrow{Separaste lo que compraste por gusto de lo que compraste para no quedar por fuera.}

\begin{infoband}{milcVerde}


\textbf{En tu cuaderno (Actividad 1 · IDENTIFICA).} Título: «Actividad 1. Lo que hay que tener». \textbf{Formato:} la lista con precios y la suma + qué tienes tú y por qué. \textbf{Extensión:} media página. \textbf{Sabes que terminaste cuando} tienes \textbf{el número total} escrito. Esta página es tuya: \textbf{nadie más tiene que leerla si tú no quieres}.
\end{infoband}

\puente{Ya la tienes. Ahora, por qué eso funciona como señal, por qué la señal se devalúa y qué dicen los estudios sobre comprar para sentirse alguien.}

\milcestacion{milcTurquesa}{Fase 2 · Sistematización: entender la señal}

\begin{softbox}{milcTurquesa}{milcGris}{Cuatro claves sobre lo que se compra}
Cuatro claves para entender por qué una marca puede costar tanto y durar tan poco.
\end{softbox}

\milcpaso{1}{milcTurquesa}{\textbf{No estás comprando la cosa: estás comprando una señal.} Unos tenis sirven para caminar, y para caminar sirven casi todos. \emph{Lo que se paga de más no es el caucho: es lo que los tenis \textbf{dicen} de ti} ---que estás al día, que en tu casa se puede, que perteneces---. \textbf{Y ahí está el detalle que lo explica todo:} \emph{una señal solo funciona si los demás la reconocen}. \textbf{Es decir: el valor no está en el objeto ni en ti ---está en la mirada de los otros---,} y por eso lo que compras nunca alcanza a ser tuyo del todo. \emph{Quien decide cuánto vale tu señal es el mismo grupo que te la exigió.}}
\milcpaso{2}{milcTurquesa}{\textbf{Toda señal comprada se devalúa, y rápido.} Cuando pocos tienen algo, tenerlo distingue. \textbf{Cuando lo tienen todos, ya no dice nada} ---y hay que comprar lo siguiente---. \emph{Por eso la lista de «lo que hay que tener» cambia cada año y nunca se acaba:} \textbf{no está diseñada para acabarse}. \emph{Y aquí conviene acordarse del momento 1:} \textbf{el peaje sube}. \textbf{La versión en plata de ese peaje es exactamente esta,} con un agravante: \emph{el que sube el precio no es una persona a la que uno pueda reclamarle}. \textbf{Es el grupo entero, y el grupo entero no responde por nada.}}
\milcpaso{3}{milcTurquesa}{\textbf{Comprar no llena el hueco que empujó la compra.} El estudio de Dittmar y su equipo, con \textbf{753 resultados de 259 estudios}, encontró que a mayor materialismo, \textbf{menor bienestar}; y la explicación que proponen es que \textbf{las necesidades psicológicas quedan mal satisfechas} (Dittmar, Bond, Hurst y Kasser, 2014). \emph{Traducido:} \textbf{uno compra buscando sentirse aceptado, y la compra no entrega aceptación} ---entrega un objeto---. \textbf{Entonces el alivio dura poco y vuelve la necesidad.} \emph{Y hay un dato que le quita culpa a la persona y se la pone al ambiente:} el efecto es más fuerte \textbf{donde el entorno apoya esos valores}. \textbf{No es que seas superficial: es que estás ocho horas al día en un lugar que califica.}}
\milcpaso{4}{milcTurquesa}{\textbf{Lo que se hace no se devalúa, porque no lo puede tener otro.} \emph{Un calado se reconoce: quien sabe mirar distingue la mano que lo hizo.} \textbf{Eso es lo contrario de una marca ---que la tienen millones y no dice nada de ninguno---.} \emph{Y funciona igual con cualquier cosa que sepas hacer:} arreglar bicicletas, cantar, cocinar, editar video, explicar matemáticas, jugar bien, dibujar, cuidar animales. \textbf{Esas señales tienen tres ventajas enormes:} \emph{no se compran, no se copian y no las devalúa que otro las tenga} ---si dos saben dibujar, ninguno vale menos---. \textbf{Y hay una cuarta, la mejor:} \emph{te dan un lugar en el grupo que no depende de la plata de tu casa}.}

\begin{drawbox}{milcTurquesa}{Las tres preguntas antes de comprar una señal}
\begin{pasoslista}
\item \textbf{¿Lo querría si nadie más lo viera?} \emph{Si la respuesta es no, no estás comprando la cosa: estás comprando la señal.}
\item \textbf{¿Cuánto va a durar la señal?} \emph{Piensa cuánto duró la anterior, la que ya nadie menciona.}
\item \textbf{¿Qué se deja de hacer en mi casa para esto?} \emph{Esta es la pregunta que casi nadie hace, y la única que se paga de verdad.}
\item \textbf{Y una cuarta, para el otro lado:} \emph{¿a quién estoy dejando por fuera cuando repito que «hay que tener» algo?}
\end{pasoslista}
\end{drawbox}

\begin{softbox}{milcMostaza}{milcGris}{Errores comunes y su corrección}
\textbf{«A mí eso no me importa»:} casi todo el mundo lo dice y casi nadie sale a la calle sin pensar qué se pone. \textbf{Creer que es un problema de vanidad:} el estudio dice que es de necesidades mal satisfechas. \textbf{Juzgar al que sí compra:} casi siempre está pagando por entrar, no presumiendo. \textbf{Juzgar al que no puede:} es el daño principal de toda esta lógica. \textbf{«Es que mis papás pueden»:} que se pueda no vuelve la señal más duradera. \textbf{Pensar que la solución es no tener nada:} la solución es tener \textbf{otra} señal, no ninguna. \textbf{Comentar la ropa de alguien:} eso es lo que hace la lista.
\end{softbox}

\begin{infoband}{milcTurquesa}


\textbf{En tu cuaderno (Actividad 2 · EXPLICA).} Título: «Actividad 2. Señal comprada, señal hecha». \textbf{Formato:} dos columnas. En la primera, tres señales compradas de tu curso, con \textbf{cuánto duraron} las del año pasado. En la segunda, tres señales \textbf{hechas} que existan en tu salón ---alguien que canta, que dibuja, que arregla cosas---. \textbf{Extensión:} media página. \textbf{Sabes que terminaste cuando} puedes explicar por qué \textbf{una señal comprada se devalúa y una hecha no}.
\end{infoband}

\puente{Ya sabes cómo opera. Es hora de lo único que cambia algo: \textbf{tener una señal que no se compre} ---algo que sepas hacer---.}

\milcestacion{milcMagenta}{Fase 3 · Praxis: hacer en vez de comprar}

\begin{softbox}{milcMagenta}{milcGris}{Cómo se cambia una señal comprada por una hecha}
Esta guía no termina en «no compres». Termina en algo más difícil y más útil: \textbf{tener otra cosa que mostrar}.
\end{softbox}

\milcpaso{1}{milcMagenta}{\textbf{El inventario y su precio (7 min).} Termina la suma. \emph{Y agrégale una columna que cambia la conversación:} \textbf{cuánto duró como señal} cada cosa parecida del año pasado. \emph{Casi siempre la respuesta está entre tres meses y un año, y verlo escrito hace más que cualquier sermón.}}
\milcpaso{2}{milcMagenta}{\textbf{Quién queda por fuera (8 min).} \textbf{Sin nombres.} Calcula \textbf{cuánta gente de tu curso no puede pagar esa lista}. \emph{No hace falta saberlo con exactitud: basta con darse cuenta de que existe.} \textbf{Y responde una pregunta:} \emph{esas personas, ¿están en los grupos, o están en la lista del momento 9 del grado 7?} \textbf{Casi siempre hay cruce, y ese cruce es el verdadero costo de la lista.}}
\milcpaso{3}{milcMagenta}{\textbf{Lo que sé hacer (10 min).} Escribe \textbf{tres cosas que sabes hacer} ---con las manos, con la cabeza, con la voz, con el cuerpo---. \emph{Si te cuesta, pregúntale a un compañero para qué te busca la gente:} \textbf{casi siempre uno es el último en saber en qué es bueno.} \textbf{Escoge una} y escribe qué te falta para hacerla mejor.}
\milcpaso{4}{milcMagenta}{\textbf{El plan de cuatro semanas (10 min).} Una señal hecha no aparece sola: \textbf{hay que volverla visible}. Escribe \textbf{qué vas a hacer cada semana durante un mes} para que eso que sabes hacer \textbf{exista para los demás}: enseñárselo a alguien, ofrecerlo para un trabajo del colegio, montarlo en la izada, arreglarle algo a alguien, meterte al semillero. \emph{Una acción por semana, concreta y con día.} \textbf{Y ojo con la trampa:} \emph{esto no es para presumir. Es para que tu lugar en el grupo dependa de algo que no te puedan quitar subiendo el precio.}}

\begin{drawbox}{milcMagenta}{Antes de cerrar tu inventario}
\begin{checklista}
\item Tu lista tiene precios reales y está \textbf{sumada}.
\item Agregaste \textbf{cuánto duró} como señal lo del año pasado.
\item Calculaste, sin nombres, \textbf{a cuánta gente deja por fuera} esa lista.
\item Escribiste tres cosas que sabes hacer, aunque te haya costado encontrarlas.
\item Tu plan tiene \textbf{cuatro acciones, una por semana, con día}.
\item Entendiste que el punto no es no tener nada: es \textbf{tener otra señal}.
\end{checklista}
\end{drawbox}

\begin{softbox}{milcMagenta}{milcGris}{Plantilla de trabajo}
\textbf{Las tres preguntas.} \emph{«¿Lo querría si nadie lo viera? · ¿Cuánto va a durar la señal? · ¿Qué se deja de hacer en mi casa por esto?»} \\[1mm]
\textbf{Lo que sé hacer.} \emph{«Sé \_\_\_\_. Para hacerlo mejor me falta \_\_\_\_.»} \\[1mm]
\textbf{Mi plan.} \emph{Semana 1: \_\_\_\_ · Semana 2: \_\_\_\_ · Semana 3: \_\_\_\_ · Semana 4: \_\_\_\_} (con día cada una).
\end{softbox}

\begin{infoband}{milcMagenta}


\textbf{En tu cuaderno (Actividad 3 · CREA).} Título: «Actividad 3. Mi inventario y mi plan». \textbf{Formato:} la lista con precios, suma y duración + el cálculo de quién queda por fuera + las tres cosas que sabes hacer + el plan de cuatro semanas. \textbf{Extensión:} una página. \textbf{Sabes que terminaste cuando} tu plan tiene \textbf{cuatro días concretos} escritos.
\end{infoband}

\puente{El producto es un inventario con precios y un plan. \emph{No es una guía contra la ropa: es una guía sobre quién decide lo que vales.}}

\milcestacion{milcMostaza}{Producto de la sesión}
\textbf{Inventario de señales: la lista con su precio real, a quién deja por fuera y un plan de cuatro semanas}.
\begin{softbox}{milcMostaza}{milcGris}{Criterios de calidad}
(1) El inventario tiene precios reales, está sumado e incluye cuánto duró la señal anterior. (2) Distingue lo comprado por gusto de lo comprado para no quedar por fuera, sin autoengaño. (3) Estima a quién excluye la lista y lo cruza con quién queda por fuera socialmente. (4) Identifica al menos una capacidad propia y qué falta para desarrollarla. (5) El plan de cuatro semanas tiene acciones concretas con día, orientadas a hacer visible esa capacidad.
\end{softbox}

\puente{Antes de cerrar, mira la parte que incomoda: en tu lista de «hay que tener», ¿cuánto suma? ¿Y quién en tu curso no puede pagarlo?}

\milcestacion{milcVino}{Evaluación liberadora · cinco dimensiones}

\begin{softbox}{milcVino}{milcGris}{Sentido de la evaluación}
Evaluar no es solo revisar una nota. Es reconocer qué aprendí, qué cambió en mí, qué emociones manejé, cómo me relaciono con la ciudad y la comunidad, y qué le dejaré a quien viene detrás.
\end{softbox}

\textbf{Mi reflexión liberadora en cinco dimensiones.} Completa cada frase iniciada:

\small
\begin{tabularx}{\textwidth}{>{\bfseries\RaggedRight\arraybackslash}p{3.4cm}Y>{\RaggedRight\arraybackslash}p{4.6cm}}
\rowcolor{milcVino}\color{white}Dimensión & \color{white}\bfseries Mi respuesta (frame starter) & \color{white}\bfseries Algo que descubrí\\
1. Personal & Lo que más cambió en mí fue \linead{6.5cm} & \linead{4.2cm}\\
2. Emocional & La emoción más fuerte fue \linead{2.5cm} y la regulé \linead{3cm} & \linead{4.2cm}\\
3. Ciudadana & En democracia esto importa porque \linead{6cm} & \linead{4.2cm}\\
4. Local (Cartago) & En mi barrio sería útil para \linead{6.5cm} & \linead{4.2cm}\\
5. Intergeneracional & A mi abuela/abuelo le diría \linead{6.5cm} & \linead{4.2cm}\\
\end{tabularx}
\normalsize

\begin{infoband}{milcVino}
\textbf{Para la próxima sustentación.} Vuelve a esta tabla en seis meses. Verás que las respuestas cambian — no porque lo aprendido cambie, sino porque tú cambias. La evaluación liberadora es una conversación contigo a través del tiempo.
\end{infoband}


\puente{Cerramos con tres voces: el que siempre es más que su apariencia, el cambio permanente de las cosas, y los pequeños datos con que se arma un juicio.}

\milcestacion{milcGradeDark}{Triángulo de pensamiento · cierre filosófico}

\begin{softbox}{milcVino}{milcGris}{Dussel · lente del nosotros}
\textit{«Analéctico quiere indicar el hecho real humano por el que todo hombre, todo grupo o pueblo se sitúa siempre más allá (aná-) del horizonte de la totalidad.»}

{\footnotesize\itshape\color{milcNegro!70}--- Enrique Dussel · Filosofía de la liberación (1977), §2.4}

Una lista de «lo que hay que tener» funciona como una \textbf{totalidad}: pretende decir, con unos tenis, todo lo que una persona vale. \emph{Dussel dice que eso es imposible ---nadie cabe ahí---.} \textbf{Y hay una consecuencia que se siente en el cuerpo:} \emph{el que entra a un grupo por lo que tiene sabe, en el fondo, que lo aceptaron por eso}, y por eso vive con miedo de que se le acabe. \textbf{Fíjate en la diferencia con el calado:} \emph{una pieza dice quién la hizo}. \textbf{Una marca no dice nada de nadie, porque la tienen millones.} \emph{Lo que se compra te vuelve igual a otros; lo que haces te vuelve reconocible.}

\textbf{Si mañana no pudiera comprar nada nuevo en un año, ¿qué quedaría de lo que los demás saben de mí?}
\end{softbox}
\linead{\linewidth}\\[1mm]
\linead{\linewidth}

\begin{softbox}{milcMostaza}{milcGris}{Estoicismo · Marco Aurelio · Meditaciones VII, 47 (c. 175 d.C.) · cuidado interior}
\textit{«Contempla las evoluciones de los astros y piensa al mismo tiempo que evolucionas con ellos; y reflexiona continuamente cómo los elementos cambian unos en otros.»}

{\footnotesize\itshape\color{milcNegro!70}--- Marco Aurelio · Meditaciones VII, 47 (c. 175 d.C.)}

Marco Aurelio se obliga a mirar lo que cambia, y una moda es \textbf{el ejemplo más rápido de eso}. \emph{Haz la prueba tú mismo:} \textbf{busca una foto de hace tres o cuatro años} ---tuya, de un familiar, de cualquiera---. \emph{Lo que en esa foto era exactamente lo que había que tener, hoy da risa.} \textbf{Y eso va a pasar con lo de ahora, sin ninguna duda.} \emph{No es un argumento para no comprar nada:} \textbf{es un argumento para no pagar un precio permanente ---la plata de la casa, la tranquilidad de tus papás, quedar debiendo--- por algo que dura una temporada.}

\textbf{Lo que hoy «hay que tener» en mi curso, ¿qué era hace tres años? ¿Y quién se acuerda?}
\end{softbox}
\linead{\linewidth}\\[1mm]
\linead{\linewidth}

\begin{softbox}{milcTurquesa}{milcGris}{Floridi · lente de la infoesfera}
\textit{«Los pequeños patrones solo pueden ser significativos si se agregan correctamente, se comparan y se procesan a tiempo.»}

{\footnotesize\itshape\color{milcNegro!70}--- Luciano Floridi · Big data and their epistemological challenge (2012)}

Un curso hace exactamente lo que describe Floridi, pero mal: \textbf{toma pequeños patrones} ---los tenis, el celular, la maleta--- y los \textbf{agrega} en un juicio completo sobre alguien. \emph{Y como en el momento 4, el problema está en la palabra \textbf{correctamente}:} \textbf{esos datos no dicen lo que se les hace decir}. \emph{Unos tenis viejos pueden significar que en la casa no alcanza, o que a esa persona le da igual, o que prefirió gastar en otra cosa.} \textbf{Hay una versión digital que aprieta más:} \emph{en las redes, lo que se ve de alguien es lo que esa persona escogió mostrar}. \textbf{Te estás comparando con la mejor toma de la vida de otro.}

\textbf{¿Con quién me comparo en redes? ¿Y qué parte de su vida no aparece ahí?}
\end{softbox}
\linead{\linewidth}\\[1mm]
\linead{\linewidth}

\begin{infoband}{milcNegro}
\textbf{Nota para el docente.} Las tres son citas textuales y llevan su referencia debajo. Puedes remitir a la obra en clase.
\end{infoband}

\milcestacion{milcVino}{Mi autoevaluación · marca una casilla por fila}

{\small
\setlength{\extrarowheight}{2.2pt}
\begin{tabularx}{\textwidth}{>{\RaggedRight\arraybackslash\bfseries}p{3.0cm}YYY}
\rowcolor{milcVino}\color{white}\bfseries Criterio & \color{white}\bfseries Lo logré & \color{white}\bfseries En proceso & \color{white}\bfseries Necesito apoyo\\[.6mm]
Comprensión del tema & \milcopcion{Explico las ideas con mis palabras.} & \milcopcion{Reconozco las ideas, falta precisión.} & \milcopcion{Necesito repasar lo básico.}\\[.8mm]
\rowcolor{milcGris}Calidad del inventario & \milcopcion{Distingo lo que compro por gusto de lo que compro para que no me dejen por fuera.} & \milcopcion{Reconozco la presión de la marca, pero creo que a mí no me afecta.} & \milcopcion{Sigo pensando que lo que uno tiene es lo que uno es.}\\[.8mm]
Aplicación y producto & \milcopcion{Mi producto cumple los criterios.} & \milcopcion{Está armado, con ajustes pendientes.} & \milcopcion{Aún está en construcción.}\\[.8mm]
\rowcolor{milcGris}Reflexión 5 dimensiones & \milcopcion{Respondí cada una con honestidad.} & \milcopcion{Respondí algunas con detalle.} & \milcopcion{Mis respuestas son muy generales.}\\[.8mm]
Triángulo de pensamiento & \milcopcion{Conecté las 3 voces con mi trabajo.} & \milcopcion{Conecté al menos una.} & \milcopcion{No las relaciono con mi proceso.}\\[.8mm]
\end{tabularx}}

\vspace{3mm}
\textbf{Hoy entendí que:}\\[1mm]
\linead{\linewidth}\\[1mm]
\linead{\linewidth}

\textbf{Mi compromiso para la próxima sesión es:}\\[1mm]
\linead{\linewidth}\\[1mm]
\linead{\linewidth}

\begin{softbox}{milcMostaza}{milcGris}{Cita de despedida · Marco Aurelio}
\textit{``El obstáculo en el camino se vuelve el camino. Lo que estorba al camino se vuelve camino.''}\\[1mm]
Cada error de hoy, bien observado, se convierte en pieza del camino que pisarás mañana.
\end{softbox}

\small
\begin{tabularx}{\textwidth}{>{\bfseries}p{3.6cm}Y}
\rowcolor{milcGradeDark!12}Nombre completo: & \linead{10cm}\\
Fecha: & \linead{4cm} \quad Curso/grupo: \linead{4cm}\\
Mi firma: & \linead{9cm}\\
\end{tabularx}
\normalsize

\begin{infoband}{milcGradeDark}
\textbf{Para la próxima sesión.} Dos cosas. \textbf{Primero, haz la acción de la semana 1} ---la que escribiste con día--- y anota qué pasó. \emph{Y haz también algo más pequeño y más difícil: pasa una semana sin comentar lo que otro trae puesto. Anota cuántas veces estuviste a punto.} \textbf{Segundo, pregunta en tu casa cómo se conseguían las cosas:} averigua con alguien mayor \textbf{qué se usaba cuando tenía tu edad, cómo lo conseguía} ---si se compraba, si se cosía en la casa, si se heredaba de un hermano---, \textbf{y si alguien se burlaba de quien no lo tenía}. \textbf{Trae:} cómo salió tu acción y lo que te contaron. \emph{Vas a encontrar dos cosas: que la presión existía igual ---no es una invención de tu generación---, y que muchas de las cosas que hoy se compran hechas antes se hacían en la casa, con las manos, y por eso llevaban el nombre de alguien.}
\end{infoband}

\end{document}
